% !TEX root = ../algebra_02.tex
\section{Сопряжённый оператор}

\begin{Definition}{Сопряженный оператор}{}
	Пусть $V$ — евклидово пространство конечной размерности, $A \in \mathrm{End}(V)$ — линейный оператор.
	Оператор $B \in \mathrm{End}(V)$ называется сопряжённым к $A$, если для любых векторов $v, w \in V$ выполняется равенство:
	\[
	(Av, w) = (v, Bw)
	\] 
\end{Definition}

\begin{Proposition}{Существование и единственность сопряженного оператора}{}
	Пусть $A \in \mathrm{End}(V)$.
	Тогда:
	\begin{enumerate}
		\item Оператор $B$, сопряжённый к $A$, существует и единственен.
		\item Если $E$ — ОНБ в пространстве $V$, то в этом базисе матрица сопряжённого оператора является транспонированной матрицей исходного оператора: если $[A]_E = A$, то $[B]_E = A^T$.
	\end{enumerate}
\end{Proposition}

\begin{proof}
	Зафиксируем произвольный ОНБ $E = (e_1, \dots, e_n)$. Пусть $A = (a_{ij})$ и $B = (b_{ij})$ — матрицы операторов $A$ и $B$ в этом базисе.
	Любые векторы $v$ и $w$ можно представить через их координаты: $v = E \cdot [v]_E$, $w = E \cdot [w]_E$.
	Равенство $(Av, w) = (v, Bw)$ должно выполняться для всех векторов, в частности, для базисных:
	\[
	(Ae_i, e_j) = (e_i, Be_j) \quad \text{для всех } i, j = 1, \dots, n.
	\] 
	Вычислим левую часть, зная, что $Ae_i$ — это вектор, координатами которого является $i$-й столбец матрицы $A$: $Ae_i = a_{1i}e_1 + \dots + a_{ni}e_n$.
	\[
	(Ae_i, e_j) = (a_{1i}e_1 + \dots + a_{ni}e_n, e_j) = a_{ji}
	\] 
	(так как базис ортонормированный, все скалярные произведения $(e_k, e_j) = 0$ при $k \neq j$, и $(e_j, e_j) = 1$).
	Аналогично вычислим правую часть: $Be_j = b_{1j}e_1 + \dots + b_{nj}e_n$, тогда:
	\[
	(e_i, Be_j) = (e_i, b_{1j}e_1 + \dots + b_{nj}e_n) = b_{ij}
	\] 
	Следовательно, условие $(Av, w) = (v, Bw)$ равносильно тому, что для всех $i, j$:
	\[
	a_{ji} = b_{ij} \iff B = A^T
	\] 
	Это доказывает как единственность сопряжённого оператора (его матрица однозначно определяется как транспонированная), так и его существование (достаточно задать оператор матрицей $A^T$ в ОНБ).
\end{proof}

\begin{Statement}{Свойства сопряжённого оператора}{}
	Пусть $A, B \in \mathrm{End}(V)$, а $\lambda \in \mathbb{R}$.
	Тогда:
	\begin{enumerate}
		\item $(A+B)^* = A^* + B^*$ 
		\item $(\lambda A)^* = \lambda A^*$ 
		\item $(AB)^* = B^* A^*$ 
		\item $(A^*)^* = A$ 
	\end{enumerate}
	Все эти свойства напрямую следуют из свойств транспонированной матрицы, так как в ОНБ матрица сопряжённого оператора равна транспонированной ($[A^*]_E = [A]_E^T$).
\end{Statement}

\begin{proof}
	Докажем свойство 3. Для любых векторов $v, w \in V$:
	\[
	(v, B^* A^* w) = (Bv, A^* w) = (ABv, w)
	\] 
	По определению сопряжённого оператора $(ABv, w) = (v, (AB)^* w)$.
	Отсюда следует, что $(AB)^* = B^* A^*$.
\end{proof}
